% SPDX-License-Identifier: CC-BY-4.0
% Copyright (c) 2026 Martin Schröder <info@swedishembedded.com>

\section{A small model teaching itself from raw documentation}\label{sec:documentation}

The coding experiment measures one generation of a 27B model learning from a teacher's
or its own verified trajectories. This section closes the whole loop on a smaller scale:
a 4B model is given nothing but raw documentation and an objective, writes its own
tasks, is graded by an interpreter, is trained on every weight, and, once trained, writes
the next round's tasks itself. No stronger model takes part.

\subsection{Setup}

\paragraph{Source and objective.} The source is three chapters of the Lua~5.1 reference
manual (string manipulation and patterns, table manipulation, mathematical functions) as
the HTML the project publishes, in its own Latin-1 encoding; the objective is the
sentence ``Become an exceptional Lua 5.1 programmer''. A local HTML page is read in the
charset it declares and converted to the Markdown the sectioner reads, exactly as a
fetched page is.\cl{SL01} With \code{--structural-graph} the skill graph is read from the
sections, one skill per function or topic, without asking a planner model.\cl{SL02}

\paragraph{Tasks and their answers.} Every task is of the kind \code{predict-lua}: a short
program and the question what it prints. The model writes the program and the question;
admission sets the program below the question in a fenced block, so it is shown once and
exactly as it is run,\cl{SL03} and the answer is what the pinned Lua~5.1.5 interpreter
prints: the program must run to completion, print a bounded, non-empty output and print
the same again, and that output becomes the reference whatever the model predicted.\cl{SL04}
An answer is graded by the formal verifier, which takes an output given in a fenced
block as the answer and treats two different blocks as a hedge.\cl{SL05}

\paragraph{The loop.} Probes are written once, before any training, by the model as it
is, and split per skill into dev probes (the baseline and every round's decision) and
final probes (the closing assessment only); a skill is asked for probes again while it
has fewer than both halves need.\cl{SL06} Each round picks the skill expected to benefit
most, has the incumbent write training tasks for it and attempt each four times, keeps
the tasks it solved sometimes and the verified answers to them, trains every weight of
the incumbent on them with a replay of the earlier accepted rounds' records, and accepts
the candidate only if the target skill improved on its dev probes and no other skill
lost more than the retention bound.\cl{SL07} With \code{--self-generate} the incumbent writes
every round's tasks: the base in round one, the last accepted checkpoint after.\cl{SL08}
The control arm is the same run without it: the base writes every round's tasks.

The records of these runs - the source, the probes, every round's stage and the
measurements this section cites - are kept in the evidence dataset
\code{splinter-flywheel-evidence}, with its analysis script.

\subsection{What full fine-tuning costs on one card}

Brain's full fine-tune holds the weights and gradients of a Qwen3 model in fp32 on the
card and the AdamW moments and master weights in host memory; \code{samples/text/chat-full-finetune}
runs it on any base in the model store.\cl{SL09} \Cref{tab:fullft} gives what eight steps
cost on one GH200 at a 631-token block.\cl{SL10}

\begin{table}[t]
  \centering
  \small
  \begin{tabular}{lrrrr}
    \toprule
    Model (optimizer state) & Card (GiB) & Host (GiB) & s / 8 steps & Loss \\
    \midrule
    Qwen3-0.6B (fp32) & 7.8 & 10.9 & 62 & 1.38 $\to$ 0.91 \\
    Qwen3-1.7B (fp32) & 15.9 & 32.6 & 128 & 1.56 $\to$ 0.81 \\
    Qwen3-4B (fp32) & 33.6 & 75.2 & 251 & 1.76 $\to$ 0.91 \\
    Qwen3-8B (fp32) & 65.6 & 152.7 & 488 & 1.12 $\to$ 0.47 \\
    Qwen3-8B (int8 moments) & 65.7 & 107.5 & 459 & 1.30 $\to$ 0.83 \\
    \bottomrule
  \end{tabular}
  \caption{Full fine-tuning on one GH200: the run's peak on the card beyond what it held
  before, the process's peak resident host memory, wall time and training loss before the
  first and after the last of eight steps (80 chat records written from the Lua manual).
  The 8B fits the card; in fp32 its host state is 153~GiB, beyond what a 205~GiB job can
  spare beside a served 27B model, and int8 moments save 45~GiB of it.}
  \label{tab:fullft}
\end{table}


\subsection{What a small model needs from the loop}\label{sec:small-model-needs}

Each of the following stopped the loop on the Lua chapters before a single round could
train; each was fixed in the code the paper describes, and each is a requirement any loop
driven by a model of this size has to meet.

\paragraph{The task must show what is run.} Asked to put its program in the material and
again verbatim in the instruction, the 4B model wrapped it in other code
(\code{material = [[print(\dots)]]}), a program that prints nothing, and admission, which
checked only that the material occurs in the instruction, accepted it. The model now
writes the question alone and admission sets the material below it.\cl{SL03}

\paragraph{The answer must come from the runtime.} With the model's own prediction as
the reference, a task was refused whenever the model mispredicted its program, and most
skills were left with too few probes. The reference is now what the interpreter
prints.\cl{SL04}

\paragraph{The verifier must read the answer the model gives.} The model answers a
predict task with a sentence above a fenced block. The verifier took candidates only
from emphasis and numbers, and most failing baseline probes held the exact output in a
fence: the baseline measured the answer's format, and training would have taught it.\cl{SL05}

\paragraph{One request is not enough.} A 4B generator admits two or three tasks per
request; probes and training tasks are asked for again, shown what was already written,
while a skill is short.\cl{SL06,SL11}

\paragraph{Generation must be compared by what it shows.} Every predict task asks the
same question of a different program; compared by its whole instruction, every pair was a
near duplicate and nearly every training task was refused as repeating a probe.
A task that shows material is now compared by the material.\cl{SL12}

\paragraph{Writing tasks is not executing a procedure.} The generator sampled at the
agent's temperature (0.2), which keeps a small model on task while it acts but has it
write the same programs whenever it is asked again; it now samples at 1.2 with a top-k of
20.\cl{SL13}

\paragraph{The environment can teach.} A task the policy never solves was taught by the
policy shown the task's passage, and a model cannot run a program; with
\code{--reference-teacher} the runtime's output is said as the teacher's reply and graded
like any answer.\cl{SL14}

\paragraph{The exam is not material.} A round's request showed the generator its skill's
probes as tasks not to repeat; the 4B model copied them, and most training tasks of a
round were exact probes. The generator is now never shown the exam.\cl{SL15}

\paragraph{An answer must be one the language defines.} A seeded \code{math.random}
program prints the same on every run and passed every check, but what it prints is the C
library generator's draw. A candidate fully fine-tuned on such answers lifted its target
skill from 0 to 4 of 4 dev probes while the other skills fell from 35 to 20 passed probes
(\code{string.byte(s, 1)} answered 97 for every string); the retention bound rejected
it.\cl{SL18}
Programs that read a generator, the clock, the environment, input or the collector are
now refused.\cl{SL16}

\paragraph{A probe must be answered greedily.} Probes were answered on the model's own
sampling (0.2, a seed per request), so the baseline and every round's decision were draws.
They are now answered greedily, and the same base on the same exam scores the same in both
arms, 34 of 56 dev probes.\cl{SL17,SL19}

\paragraph{The loop learns at a small step.} At a learning rate of $10^{-6}$ (six steps,
one verified answer per task, retention bound 0.5) rounds are accepted: each arm accepted 5
rounds, and on the final probes, which no round trains on, the self-generate arm went from
26 to 35 of 53 passes (7 skills up, 1 down) and the fixed-generator arm from 26 to 40 of 53
(8 up, 0 down).\cl{SL20} The control did better than the arm whose trained model writes the
next tasks, so these runs show that the loop acquires skills from raw documentation, not
that self-generated tasks help.

\paragraph{A second replicate.} The same experiment on a fresh exam, written by the base model,
with the full fine-tune of a later brain revision, accepted 3 rounds in each arm. On the final probes the self-generate arm went
from 20 to 30 of 46 in 13 rounds (7 skills up, 0 down) and the fixed-generator arm from 20 to 29 of
46 in 8 rounds (7 up, 2 down).\cl{SL21} Both arms learn in both replicates; which generator does
better changes sign between them, so on this evidence the loop learns from tasks its own model
writes, and whether the trained model's tasks are better than its base's is not settled. A third
replicate on the second exam, which changes only the runs' random draws, agrees: the self-generate arm
went from 20 to 33 of 46 (2 of 10 rounds accepted, 8 skills up, 0 down) and the fixed-generator arm
from 20 to 32 (3 of 9 accepted, 7 up, 1 down).\cl{SL22} All six arms of the three replicates ended
above their baseline on probes no round trained on; the self-generate arm finished ahead in two
replicates and behind in one.
